\documentclass[IB]{TFUOC}%IB: CASTELLANO, CAT: CATALÁN, ENG: ANGLÈS

\usepackage{float}
\usepackage{pgfgantt}
\usepackage{tabularx}
\usepackage{booktabs}
\usepackage{array}

\setlength{\parskip}{0.6em}
\setlength{\parindent}{0pt}

%Introducción de datos del trabajo
\title{Predicción de precios de vivienda en Andalucía mediante aprendizaje automático}
\titcrt{Predicción de precios de vivienda en Andalucía}

\author{Mario Aroca Páez}
\date{25 de octubre de 2026}


\nomPDC{Aarón Montero Montero}
\nomPRA{Susana Acedo}
\titulac{Máster en Ciencia de Datos}
\area{Trabajo Final de Máster}
\idioma{Castellano}
\credits{15}
\parcla{mercado inmobiliario, ciencia de datos, aprendizaje automático, predicción de precios, Andalucía}

\licenc{ccBy}
%Posibles licencias
%ccByNcNd
%ccByNcSa
%ccByNc
%ccByNd
%ccBySa
%ccBy
%GNU
%copyright


%%%%%%%%%%
% Resumen en el idioma

\abstractidioma{
Este Trabajo Final de Máster tiene como objetivo analizar el mercado inmobiliario residencial en Andalucía mediante técnicas de ciencia de datos y aprendizaje automático. El trabajo se centrará en la construcción de modelos capaces de estimar el precio de oferta de una vivienda a partir de sus características y de información relacionada con su localización y entorno.

Para ello, se llevará a cabo un proceso completo de análisis de datos que incluirá la obtención y preparación de la información, el análisis exploratorio, el tratamiento de valores ausentes y atípicos, la transformación de variables y la construcción de distintos modelos predictivos. Se plantea comparar varios enfoques de aprendizaje automático con un modelo base, utilizando métricas de evaluación adecuadas para determinar su rendimiento.

Además de buscar una buena capacidad predictiva, el trabajo prestará especial atención a la interpretabilidad de los modelos. Para ello se utilizarán técnicas como SHAP, con el objetivo de identificar qué variables tienen una mayor influencia en la estimación del precio.

El proyecto también tendrá en cuenta las limitaciones de los datos, los posibles sesgos derivados de las fuentes utilizadas y los aspectos éticos relacionados con su obtención y tratamiento.
}

% Resumen en inglés.
\abstractenglish{
This Master's Thesis aims to analyse the residential real estate market in Andalusia using data science and machine learning techniques. The project will focus on building models capable of estimating the asking price of a property based on its characteristics and on information related to its location and surrounding area.

To achieve this, a complete data analysis process will be carried out, including data collection and preparation, exploratory data analysis, treatment of missing values and outliers, variable transformation and the development of different predictive models. Several machine learning approaches will be compared against a baseline model using suitable evaluation metrics.

In addition to predictive performance, special attention will be paid to model interpretability. Techniques such as SHAP will be used to identify which variables have the greatest influence on price estimations.

The project will also consider the limitations of the available data, possible biases derived from the selected sources and the ethical aspects related to data collection and processing.
}

\begin{document}

\estructura

\tableofcontents

%\listoffigures

%\listoftables




\chapter{Introducción}

\section{Contexto y justificación del trabajo}

El mercado de la vivienda tiene una relevancia importante desde el punto de vista económico y social y presenta una evolución condicionada por factores como la demanda, la oferta, la financiación y las características territoriales \cite{bde_vivienda2024}. A nivel individual, el precio de los inmuebles puede variar de forma considerable en función de factores como la localización, la superficie, el número de habitaciones, el estado de la vivienda o las características de la zona en la que se encuentra.

En los últimos años, la disponibilidad de grandes volúmenes de datos inmobiliarios ha permitido aplicar técnicas de ciencia de datos y aprendizaje automático para analizar este tipo de información y construir modelos de predicción de precios \cite{houseprice_survey2023}. Estos métodos permiten detectar relaciones entre variables que pueden resultar difíciles de observar mediante técnicas más tradicionales.

El presente trabajo se centrará en el mercado inmobiliario residencial de Andalucía. El objetivo será estudiar qué variables están más relacionadas con el precio de oferta de una vivienda y analizar hasta qué punto es posible estimarlo mediante diferentes modelos de aprendizaje automático.

La relevancia del trabajo no se limita únicamente a obtener una predicción precisa. También se pretende analizar el comportamiento de los modelos y entender qué factores influyen en sus resultados. Por este motivo, se incorporarán técnicas de interpretabilidad que permitan estudiar el peso de las distintas variables en las predicciones realizadas.

De esta forma, el trabajo busca combinar una parte predictiva con una parte de análisis e interpretación, siguiendo un proceso completo de ciencia de datos desde la preparación inicial de la información hasta la evaluación de los modelos y la obtención de conclusiones.

\section{Motivación}

La elección de este tema parte de mi interés por aplicar las técnicas aprendidas durante el Máster en Ciencia de Datos a un problema real y fácilmente interpretable. El mercado inmobiliario resulta especialmente adecuado para este propósito, ya que combina variables numéricas, categóricas, geográficas y socioeconómicas, lo que permite abordar distintas etapas de un proyecto completo de ciencia de datos.

Además, se trata de un tema especialmente cercano para una parte importante de la población joven que se encuentra actualmente en proceso de emancipación o de acceso a una vivienda. Las dificultades relacionadas con los precios de compra y alquiler, la disponibilidad de vivienda y las diferencias entre zonas hacen que el mercado residencial tenga una relevancia directa en decisiones personales y económicas. Este contexto aumenta mi interés por estudiar el problema desde una perspectiva basada en datos y comprender mejor qué factores están relacionados con el precio de la vivienda.

También me interesa especialmente la posibilidad de no limitar el trabajo a la obtención de un modelo predictivo, sino analizar por qué el modelo toma determinadas decisiones y qué variables tienen una mayor influencia en la estimación del precio. Esto permite trabajar no solo la precisión de los modelos, sino también su interpretación.

Por último, el ámbito geográfico de Andalucía ofrece un contexto suficientemente amplio y heterogéneo como para analizar diferencias entre zonas y estudiar el comportamiento de los modelos en distintos entornos dentro de una misma comunidad autónoma.

\section{Objetivos del trabajo}

\subsection{Objetivo principal}

El objetivo principal de este trabajo es desarrollar y evaluar modelos de aprendizaje automático capaces de estimar el precio de oferta de viviendas residenciales en Andalucía, analizando además qué variables tienen una mayor influencia en las predicciones realizadas.

\subsection{Objetivos parciales}

Para alcanzar el objetivo principal se plantean los siguientes objetivos parciales:

\begin{itemize}
\item Obtener y preparar un conjunto de datos de viviendas residenciales en Andalucía adecuado para la estimación del precio de oferta, evaluando su cobertura territorial, calidad y representatividad.
\item Realizar un análisis exploratorio de los datos para estudiar la distribución de las variables, detectar valores ausentes, valores atípicos, registros duplicados y posibles inconsistencias.
\item Preparar las variables para el entrenamiento mediante las transformaciones necesarias, incluyendo cuando proceda codificación de variables categóricas, escalado de variables numéricas, creación y selección de características y tratamiento de posibles fuentes de fuga de información.
\item Construir un modelo de regresión lineal que actúe como referencia para valorar la mejora obtenida mediante técnicas de aprendizaje automático.
\item Entrenar y comparar Random Forest y diferentes algoritmos de Gradient Boosting, concretamente XGBoost, LightGBM y CatBoost, valorando además la utilización de una red neuronal tipo MLP como modelo complementario cuando las características de los datos lo permitan.
\item Ajustar los hiperparámetros de los modelos que lo requieran mediante Optuna, utilizando exclusivamente los datos de entrenamiento y validación durante este proceso.
\item Evaluar los modelos utilizando principalmente MAE, RMSE y $R^2$, aplicando el mismo procedimiento de evaluación para garantizar una comparación homogénea entre ellos.
\item Analizar mediante SHAP la contribución de las variables a las predicciones de los modelos seleccionados, tanto desde una perspectiva global como para observaciones concretas.
\item Estudiar si la influencia de las variables y el comportamiento de los modelos presentan diferencias entre distintos ámbitos territoriales dentro de Andalucía, siempre que la cobertura de los datos permita realizar este análisis de forma representativa.
\item Identificar las principales limitaciones del estudio relacionadas con las fuentes de datos, la cobertura geográfica, la capacidad de generalización de los modelos y las decisiones metodológicas adoptadas, proponiendo posibles líneas de trabajo futuro.
\end{itemize}

\subsection{Hipótesis de trabajo}

Se plantea como hipótesis principal que los modelos de aprendizaje automático capaces de representar relaciones no lineales, especialmente los métodos basados en conjuntos de árboles y Gradient Boosting, obtendrán un rendimiento predictivo superior al modelo de regresión lineal utilizado como referencia en la estimación del precio de oferta. Asimismo, se espera que la incorporación de características propias de la vivienda junto con variables relacionadas con su localización y entorno contribuya a mejorar la capacidad predictiva. Mediante SHAP se analizará además si las variables con mayor influencia sobre las estimaciones son coherentes con los factores identificados en la literatura revisada.

\section{Impacto en sostenibilidad, ético-social y de diversidad}
\label{s:etic}

El desarrollo de este trabajo requiere prestar atención a distintos aspectos relacionados con el uso responsable de los datos y con las posibles limitaciones derivadas de las fuentes utilizadas.

Desde el punto de vista ético, se tendrá en cuenta la procedencia de los datos y las condiciones de uso de las distintas fuentes. El objetivo del trabajo no requiere tratar datos personales, por lo que se evitará almacenar información como nombres, teléfonos, correos electrónicos u otros datos que permitan identificar a una persona.

También se tendrán en cuenta los posibles sesgos presentes en los datos. La información procedente de anuncios inmobiliarios puede no representar de forma homogénea todas las zonas geográficas ni todos los tipos de vivienda, por lo que los resultados deberán interpretarse teniendo en cuenta estas limitaciones.

Desde el punto de vista de la diversidad, se tendrá en cuenta que el mercado inmobiliario andaluz presenta diferencias importantes entre provincias, municipios y tipos de entorno. Estas diferencias pueden dar lugar a comportamientos distintos en los precios y en la disponibilidad de datos, por lo que se analizará la representatividad de las distintas zonas y se evitará extraer conclusiones generales a partir de áreas o tipos de vivienda insuficientemente representados.

Desde el punto de vista social, el análisis puede contribuir a comprender mejor qué factores están relacionados con el precio de la vivienda y cómo varían entre distintas zonas. Aun así, los resultados del trabajo tendrán una finalidad analítica y académica y no se plantean como una herramienta automática para la toma de decisiones que puedan afectar directamente a personas.

En relación con los Objetivos de Desarrollo Sostenible, el trabajo se relaciona principalmente con el ODS 11, Ciudades y comunidades sostenibles, ya que incluye entre sus metas el acceso a una vivienda adecuada, segura y asequible \cite{onu_ods}. De forma secundaria, también puede relacionarse con el ODS 10, Reducción de las desigualdades, especialmente en el análisis de diferencias socioeconómicas entre zonas.


\section{Enfoque y método seguido}

El trabajo seguirá un enfoque cuantitativo y aplicado, basado en el análisis de datos y en la comparación de distintos modelos de aprendizaje automático para un problema de regresión.

Tomando como referencia las estrategias de investigación propuestas por Oates \cite{oates2006} para proyectos en el ámbito de los sistemas de información y la informática, el trabajo se aproximará principalmente a una estrategia experimental. El objetivo será evaluar de forma sistemática el comportamiento de distintos modelos predictivos utilizando los mismos conjuntos de datos y métricas de evaluación, con el fin de comparar sus resultados y analizar en qué medida permiten estimar el precio de oferta de una vivienda.

Como alternativas se han considerado otras estrategias de investigación habituales en el ámbito de los sistemas de información, como el estudio de caso o el diseño y creación de un artefacto. Sin embargo, estas opciones no se ajustan tanto al objetivo principal del trabajo, ya que no se pretende analizar en profundidad una única organización o contexto específico ni desarrollar una aplicación como producto principal. Por este motivo, se considera más adecuada una estrategia experimental, centrada en comparar de forma controlada distintos modelos predictivos y analizar su comportamiento sobre un mismo problema y conjunto de datos.

La investigación se basará principalmente en datos cuantitativos. La información procederá de anuncios inmobiliarios correspondientes al mercado residencial de Andalucía, principalmente de Fotocasa, y se estudiará su posible enriquecimiento mediante información procedente de fuentes oficiales, especialmente del Ministerio de Vivienda y Agenda Urbana (MIVAU) \cite{mivau_transacciones}. También podrán incorporarse otras fuentes oficiales cuando aporten variables relevantes para el análisis.

La estrategia elegida consiste en desarrollar un proceso completo de ciencia de datos, desde la obtención y preparación de la información hasta la evaluación e interpretación de los modelos. Este enfoque permite estudiar el problema desde varias perspectivas: calidad de los datos, comportamiento de las variables, capacidad predictiva de los modelos e interpretación de los resultados.

El proceso de trabajo se dividirá en varias fases. En primer lugar, se realizará la obtención y preparación del conjunto de datos. A continuación, se llevará a cabo un análisis exploratorio para estudiar la distribución de las variables, detectar valores ausentes, valores atípicos y posibles inconsistencias.

Posteriormente, se aplicarán las transformaciones necesarias para preparar los datos antes del entrenamiento de los modelos. Una vez completada la preparación inicial de los datos, estos se dividirán en conjuntos de entrenamiento, validación y prueba. El conjunto de prueba se mantendrá separado durante el desarrollo y se utilizará únicamente para la evaluación final. Las transformaciones que dependan de los datos y el ajuste de hiperparámetros se realizarán sin utilizar información procedente de este conjunto.

Se definirá un modelo de regresión lineal que actúe como referencia inicial de rendimiento. A partir de la revisión del estado del arte, se ha decidido centrar la comparación en Random Forest y en los algoritmos de Gradient Boosting XGBoost \cite{chen2016xgboost}, LightGBM \cite{ke2017lightgbm} y CatBoost \cite{prokhorenkova2018catboost}. Además, se valorará la utilización de una red neuronal tipo MLP como modelo complementario cuando las características y el volumen del conjunto de datos lo permitan.

Los modelos se evaluarán principalmente mediante RMSE (Root Mean Squared Error), MAE (Mean Absolute Error) y el coeficiente de determinación $R^2$. Estas métricas permitirán analizar el rendimiento desde perspectivas complementarias y se aplicarán bajo un mismo procedimiento de evaluación para facilitar la comparación entre los modelos. Cuando sea necesario, se realizará además un ajuste de hiperparámetros mediante Optuna \cite{akiba2019optuna}, utilizando únicamente los datos de entrenamiento y validación durante este proceso. Finalmente, se aplicarán técnicas de interpretabilidad, principalmente SHAP \cite{lundberg2017shap}, para analizar qué variables presentan una mayor influencia en las predicciones obtenidas.

El desarrollo se realizará principalmente mediante Python, utilizando Jupyter Notebook y Visual Studio Code como entornos de trabajo. Se emplearán librerías de análisis, visualización y aprendizaje automático adecuadas para cada fase del proyecto. La selección concreta de algunas herramientas podrá ajustarse en función de las características del conjunto de datos y de las necesidades que aparezcan durante la implementación.

Si el tiempo disponible y los resultados obtenidos lo permiten, se estudiará también la posibilidad de desarrollar una herramienta sencilla de visualización o un dashboard interactivo que facilite la exploración e interpretación de los resultados. Para ello se podrán emplear herramientas como Streamlit y Plotly. Esta parte tendrá carácter complementario y no constituye el objetivo principal del trabajo.

\section{Planificación del trabajo}

La planificación del trabajo se ha organizado tomando como referencia las distintas entregas establecidas en la asignatura. El objetivo es distribuir las tareas de forma progresiva, de manera que cada fase sirva como base para la siguiente.

La primera fase correspondió a la definición y planificación del TFM, con fecha de entrega el 4 de octubre de 2026. Durante esta etapa se concretaron el tema, la motivación, los objetivos, la hipótesis, la metodología inicial, los aspectos éticos y la planificación general del proyecto.

La segunda fase se ha centrado en la elaboración del estado del arte, cuya entrega está prevista para el 25 de octubre de 2026. Durante esta etapa se ha realizado una revisión bibliográfica sobre la predicción de precios de vivienda, el uso de técnicas de aprendizaje automático en el ámbito inmobiliario, las fuentes de datos utilizadas y los métodos de interpretabilidad aplicables al problema. Esta revisión ha permitido además concretar los modelos, métricas y procedimientos de evaluación que se utilizarán en las siguientes fases del trabajo.

La tercera fase corresponderá al diseño e implementación del trabajo y finalizará el 15 de diciembre de 2026. Esta será la etapa principal de desarrollo e incluirá la obtención y preparación de los datos, el análisis exploratorio, el tratamiento de valores ausentes y atípicos, la transformación de variables, el entrenamiento de modelos, el ajuste de hiperparámetros y la evaluación de resultados.

A continuación, se realizará la redacción y revisión de la memoria. Se prevé una entrega preliminar el 20 de diciembre de 2026 y una entrega final el 27 de diciembre de 2026. Durante esta fase se integrarán los resultados obtenidos, se revisarán las conclusiones y se incorporarán las correcciones necesarias.

La presentación audiovisual del trabajo se entregará el 5 de enero de 2027 y la documentación destinada al tribunal el 8 de enero de 2027. Finalmente, la defensa pública del TFM está prevista para el 29 de enero de 2027.

La planificación podrá ajustarse en función de las dificultades que aparezcan durante el desarrollo, especialmente en lo relacionado con la obtención y calidad de los datos, el tiempo de entrenamiento de los modelos o la necesidad de realizar cambios en la metodología.

Las distintas fases presentan una dependencia principalmente secuencial. La definición inicial del proyecto ha servido de base para el estado del arte, que ha permitido concretar las principales decisiones metodológicas antes de la fase de implementación. A su vez, la preparación y análisis de los datos deberán completarse antes del entrenamiento y evaluación de los modelos, mientras que la redacción de la memoria se irá desarrollando de forma paralela y se completará una vez obtenidos los resultados finales.

\subsection{Recursos necesarios}

Para el desarrollo del proyecto se utilizará principalmente un equipo personal con capacidad suficiente para realizar el tratamiento de los datos y el entrenamiento de los modelos planteados. El trabajo se desarrollará principalmente en Python, utilizando Jupyter Notebook y Visual Studio Code como entornos de trabajo.

Se emplearán librerías de análisis, visualización y aprendizaje automático adecuadas para cada fase del proyecto. Además, se utilizará Overleaf para la elaboración de la memoria en \LaTeX{} y Git para el control de versiones del código y de los principales recursos desarrollados.

Los datos procederán principalmente de Fotocasa y de fuentes oficiales como el Ministerio de Vivienda y Agenda Urbana. En caso de que algún modelo requiera una capacidad de cómputo superior a la disponible localmente, se valorará el uso puntual de recursos de computación en la nube.

\subsection{Cronograma de trabajo}

El cronograma de la Figura \ref{fig:cronograma} resume la planificación temporal del proyecto. Cada columna representa aproximadamente una semana de trabajo, mientras que los hitos indican las principales fechas de entrega.

\begin{figure}[H]
\centering

\begin{ganttchart}[
    hgrid,
    vgrid,
    x unit=0.47cm,
    y unit chart=0.65cm,
    y unit title=0.7cm,
    bar height=0.55,
    milestone height=0.55,
    title label font=\scriptsize\bfseries,
    bar label font=\scriptsize,
    milestone label font=\scriptsize\itshape,
    bar label node/.append style={align=right, text width=4.7cm},
    milestone label node/.append style={align=right, text width=4.7cm}
]{1}{19}

\gantttitle{Sep}{1}
\gantttitle{Oct}{4}
\gantttitle{Nov}{4}
\gantttitle{Dic}{6}
\gantttitle{Ene}{4} \\

\gantttitlelist{1,...,19}{1} \\

\ganttbar{Definición y planificación}{1}{2} \\
\ganttbar{Estado del arte}{2}{5} \\
\ganttbar{Obtención y preparación de datos}{6}{7} \\
\ganttbar{Análisis exploratorio}{8}{9} \\
\ganttbar{Preprocesamiento}{9}{10} \\
\ganttbar{Entrenamiento y ajuste}{10}{11} \\
\ganttbar{Evaluación e interpretabilidad}{11}{12} \\
\ganttbar{Redacción y revisión}{13}{14} \\
\ganttbar{Presentación audiovisual}{15}{15} \\
\ganttbar{Preparación de la defensa}{16}{19} \\

\ganttmilestone{M1 -- 04/10}{2} \\
\ganttmilestone{M2 -- 25/10}{5} \\
\ganttmilestone{M3 -- 15/12}{12} \\
\ganttmilestone{Memoria preliminar -- 20/12}{13} \\
\ganttmilestone{Memoria final -- 27/12}{14} \\
\ganttmilestone{Audiovisual -- 05/01}{15} \\
\ganttmilestone{Tribunal -- 08/01}{16} \\
\ganttmilestone{Defensa -- 29/01}{19}

\end{ganttchart}

\caption{Cronograma inicial del Trabajo Final de Máster.}
\label{fig:cronograma}
\end{figure}

\subsection{Riesgos y medidas de mitigación}

Durante el desarrollo del proyecto pueden aparecer distintos riesgos que afecten a la planificación prevista. Uno de los principales riesgos está relacionado con la disponibilidad y calidad de los datos. Es posible encontrar valores ausentes, registros duplicados, variables poco fiables o una representación desigual entre distintas zonas geográficas. Para reducir este riesgo, se realizará una revisión inicial de calidad de los datos y, cuando sea necesario, se valorará el uso de fuentes alternativas o la eliminación de variables que no ofrezcan suficiente fiabilidad.

Otro posible riesgo es que la obtención de los datos requiera más tiempo del previsto o que alguna de las fuentes seleccionadas no pueda utilizarse finalmente. Como medida de mitigación, se priorizará el uso de fuentes estables y se mantendrán alternativas que permitan continuar el trabajo sin depender de una única fuente de información.

También puede producirse un aumento del tiempo de entrenamiento de algunos modelos, especialmente durante el ajuste de hiperparámetros. En este caso, se limitará el espacio de búsqueda, se priorizarán los modelos más relevantes y se adaptará la complejidad de los experimentos a los recursos computacionales disponibles.

Por último, existe el riesgo de que alguna de las técnicas inicialmente planteadas no proporcione resultados satisfactorios. La metodología se ha planteado de forma flexible para permitir sustituir o descartar modelos cuando los resultados obtenidos o la revisión del estado del arte indiquen que existen alternativas más adecuadas.

\section{Declaración sobre el uso de inteligencia artificial generativa}

Durante el desarrollo de este Trabajo Final de Máster se utilizarán herramientas de inteligencia artificial generativa como apoyo en distintas fases del proyecto. En concreto, se prevé utilizar ChatGPT, Claude y Gemini como herramientas auxiliares para la revisión y mejora de código, la comprobación de actividades y entregables, la resolución de dudas técnicas, el apoyo en la redacción y la asistencia en la elaboración del documento en \LaTeX.

Estas herramientas podrán emplearse a lo largo de las distintas etapas del trabajo, especialmente durante la preparación de los datos, el desarrollo y revisión del código, la interpretación de errores, la estructuración de la memoria y la revisión de los distintos apartados antes de su entrega.

El uso de inteligencia artificial generativa tendrá siempre un carácter de apoyo. Las respuestas obtenidas serán revisadas y contrastadas antes de incorporarse al trabajo, especialmente en aquellas cuestiones relacionadas con decisiones metodológicas, implementación técnica, interpretación de resultados o referencias bibliográficas.

El principal beneficio esperado de estas herramientas es agilizar tareas de revisión, detección de errores y exploración de alternativas. No obstante, también se tendrán en cuenta sus limitaciones, como la posibilidad de generar respuestas incorrectas, incompletas o poco adecuadas al contexto del proyecto. Por este motivo, no se considerará ninguna respuesta generada como válida de forma automática.

El uso de estas herramientas ha sido comentado con el director del TFM como parte del apoyo previsto durante el desarrollo del proyecto. En cualquier caso, las decisiones metodológicas, el código finalmente utilizado, el análisis de los datos, la interpretación de los resultados, las conclusiones y el contenido final de la memoria serán responsabilidad del autor.


\chapter{Estado del arte}

\section{Metodología de la revisión bibliográfica}

La elaboración del estado del arte se ha basado en una revisión documental de literatura científica relacionada con la predicción de precios inmobiliarios, los modelos de aprendizaje automático aplicados a este problema y las técnicas utilizadas para interpretar sus resultados.

La búsqueda de referencias se realizó principalmente mediante Google Scholar, Web of Science, Semantic Scholar, Dialnet, ScienceDirect y el buscador de la Biblioteca de la UOC. También se utilizó Claude como herramienta de apoyo para identificar posibles trabajos relacionados y ampliar los términos de búsqueda. Las referencias localizadas mediante herramientas de inteligencia artificial fueron posteriormente comprobadas mediante la publicación original, el DOI o repositorios académicos antes de incorporarlas al trabajo.

Las búsquedas se realizaron utilizando combinaciones de términos como \textit{house price prediction}, \textit{housing price prediction}, \textit{machine learning real estate}, \textit{hedonic housing price}, \textit{explainable artificial intelligence real estate}, \textit{SHAP housing price}, \textit{asking price prediction} y sus equivalentes en castellano. También se realizaron búsquedas específicas relacionadas con España y Andalucía para localizar estudios desarrollados en mercados inmobiliarios próximos al ámbito geográfico de este trabajo.

Para los estudios aplicados se dio prioridad a publicaciones recientes, especialmente trabajos publicados entre 2020 y 2026, con el objetivo de recoger metodologías y algoritmos representativos del estado actual de la predicción inmobiliaria mediante aprendizaje automático. Este criterio temporal no se aplicó a las referencias fundamentales necesarias para justificar métodos consolidados, por lo que también se incluyeron trabajos anteriores como los de Rosen sobre modelos hedónicos, Breiman sobre Random Forest o las publicaciones originales de los principales algoritmos utilizados.

Como criterios de inclusión se consideraron principalmente artículos científicos revisados por pares relacionados con la predicción o valoración de vivienda mediante métodos estadísticos o de aprendizaje automático, trabajos que describieran las fuentes de datos, las variables, los modelos o las métricas utilizadas, y estudios que incorporasen técnicas de interpretabilidad relevantes para el objetivo del TFM. Se prestó especial atención a investigaciones desarrolladas en España y Andalucía por su proximidad al mercado analizado. También se incluyeron publicaciones originales de los algoritmos considerados y fuentes institucionales cuando resultaban necesarias para contextualizar el mercado inmobiliario o las fuentes de datos.

Se descartaron trabajos cuya relación con la predicción o valoración residencial fuera únicamente indirecta, publicaciones sin información suficiente sobre la metodología empleada y estudios centrados en problemas inmobiliarios distintos del objetivo del presente trabajo cuando no aportaban elementos metodológicos relevantes. También se evitó incorporar referencias únicamente por similitud temática cuando no contribuían a justificar alguna decisión relacionada con los datos, los modelos, la evaluación o la interpretabilidad.

La selección final de referencias se realizó atendiendo a su relación con las decisiones metodológicas del TFM y no únicamente a su fecha de publicación. De este modo, la revisión combina trabajos fundacionales con estudios aplicados recientes y antecedentes específicamente desarrollados en España y Andalucía. La revisión no pretende constituir una revisión sistemática exhaustiva de toda la literatura existente, sino identificar los principales enfoques, fuentes de datos, modelos y técnicas de evaluación e interpretabilidad necesarios para fundamentar el trabajo.

\section{Predicción de precios en el mercado inmobiliario}

La predicción del precio de la vivienda es un problema que ha sido abordado mediante diferentes técnicas estadísticas y, más recientemente, mediante modelos de aprendizaje automático. El precio de un inmueble depende de múltiples factores relacionados tanto con las características propias de la vivienda como con su localización y el entorno en el que se encuentra. Entre estos factores pueden encontrarse la superficie, el número de habitaciones, la antigüedad del inmueble, su estado de conservación o determinadas características de la zona.

Esta combinación de variables hace que la estimación del precio sea un problema complejo, ya que las relaciones entre las características de una vivienda y su precio no tienen por qué ser lineales y pueden variar entre diferentes mercados inmobiliarios y áreas geográficas. En este contexto, el aprendizaje automático permite trabajar con relaciones más complejas entre las variables y comparar diferentes enfoques predictivos.

La revisión realizada por Geerts et al. \cite{houseprice_survey2023} analiza numerosos trabajos relacionados con la predicción de precios de vivienda y muestra la variedad de fuentes de datos y metodologías utilizadas en este ámbito. La literatura incluye desde modelos estadísticos tradicionales hasta técnicas de aprendizaje automático capaces de trabajar con conjuntos de datos de mayor tamaño y con relaciones no lineales entre las variables.

En el contexto español, Mora-Garcia et al. \cite{moragarcia2022housing} aplicaron diferentes algoritmos de aprendizaje automático a datos georreferenciados del mercado inmobiliario de Alicante. El estudio compara modelos basados en boosting y bagging con una regresión lineal y muestra que los modelos de aprendizaje automático pueden adaptarse mejor a las relaciones no lineales presentes en este tipo de datos. Este trabajo resulta especialmente relevante para el presente TFM por desarrollarse sobre un mercado inmobiliario español y por utilizar varios de los algoritmos considerados posteriormente en este trabajo.

La componente geográfica también desempeña un papel importante en la valoración inmobiliaria. Guliker et al. \cite{guliker2022spatial} compararon distintos modelos de valoración incorporando información espacial, entre ellos regresión lineal, regresión geográficamente ponderada y XGBoost. Sus resultados muestran la utilidad de incorporar variables relacionadas con la localización y refuerzan la importancia de tener en cuenta las diferencias territoriales en los modelos de predicción de precios.

\subsection{Diferencias entre precio de oferta, precio de transacción y valor de tasación}

Al comparar trabajos de predicción inmobiliaria es importante distinguir entre los diferentes valores que pueden utilizarse como variable objetivo. El precio de oferta corresponde al importe solicitado por el vendedor cuando una vivienda se publica en el mercado, mientras que el precio de transacción representa el importe finalmente acordado cuando se realiza la compraventa. Ambos valores pueden diferir como consecuencia del proceso de negociación y de las condiciones concretas de la operación.

Por otra parte, el valor de tasación es una estimación del valor del inmueble obtenida mediante un procedimiento de valoración y no tiene por qué coincidir exactamente con el precio de oferta ni con el precio finalmente pagado. Por este motivo, los resultados obtenidos en estudios que utilizan valores de tasación, precios de transacción o precios procedentes de anuncios no son completamente equivalentes y deben interpretarse teniendo en cuenta la naturaleza de la variable objetivo utilizada.

Esta diferencia es especialmente relevante para el presente trabajo, cuyo objetivo es estimar el precio de oferta de viviendas residenciales. Los estudios basados en precios de transacción o valores de tasación resultan útiles para analizar metodologías, variables y algoritmos de predicción, pero sus resultados no pueden trasladarse directamente al problema planteado sin tener en cuenta estas diferencias.

\section{Fuentes de datos y variables utilizadas}

Las fuentes de datos utilizadas en la predicción de precios de vivienda pueden ser muy diversas. Además de las características propias del inmueble, muchos estudios incorporan información relacionada con la localización, el entorno urbano o las características socioeconómicas de la zona. La combinación de distintas fuentes permite construir una representación más completa de cada vivienda y analizar en qué medida estos factores contribuyen a explicar las diferencias de precio.

Geerts et al. \cite{houseprice_survey2023} muestran en su revisión que los trabajos de predicción inmobiliaria utilizan diferentes tipos de información. Entre las variables más habituales se encuentran las características estructurales de la vivienda, como la superficie, el número de habitaciones o la antigüedad, junto con variables espaciales y características relacionadas con el entorno. La disponibilidad de nuevas fuentes de información ha permitido ampliar progresivamente este conjunto de variables y desarrollar modelos capaces de aprovechar datos de distinta naturaleza.

Un ejemplo especialmente próximo al contexto de este trabajo es el estudio de Mora-Garcia et al. \cite{moragarcia2022housing}, en el que se utilizan microdatos georreferenciados procedentes del mercado inmobiliario de Alicante y se combinan con información catastral, indicadores sociodemográficos y económicos y datos obtenidos mediante imágenes de satélite. Este enfoque permite incorporar al modelo tanto características propias de las viviendas como información relacionada con el entorno en el que se encuentran.

La importancia de las variables espaciales también se observa en el trabajo de Guliker et al. \cite{guliker2022spatial}. En este estudio, variables como la superficie habitable, determinadas medidas relacionadas con la valoración de viviendas próximas, la localización geográfica, el nivel de renta de la zona o la distancia a determinados servicios contribuyen a explicar las diferencias en los valores inmobiliarios.

La información contenida en los propios anuncios también puede aportar características adicionales para los modelos. Baur et al. \cite{baur2023automated} estudian la valoración inmobiliaria mediante modelos de aprendizaje automático incorporando, además de variables estructuradas, representaciones de las descripciones textuales de los inmuebles. Sus resultados muestran que la información presente en los anuncios puede complementar a las características estructurales tradicionales y ponen de manifiesto la variedad de información que puede aprovecharse en este tipo de problemas.

Estos trabajos muestran que la calidad de un modelo de predicción no depende únicamente del algoritmo utilizado, sino también de la información disponible y de la capacidad para representar adecuadamente las características de la vivienda y de su entorno. Por este motivo, en el presente TFM se plantea utilizar información procedente de anuncios inmobiliarios y estudiar su posible enriquecimiento mediante fuentes oficiales que aporten variables adicionales relacionadas con la localización y el contexto territorial.

\subsection{Obtención y generación de los datos del TFM}

La fuente principal de información del presente trabajo estará formada por anuncios de viviendas publicados en Fotocasa. A diferencia de fuentes estadísticas agregadas, los anuncios permiten trabajar a nivel de inmueble y disponer de información como el precio de oferta, la superficie, el número de habitaciones, la localización y otras características de la vivienda que puedan encontrarse disponibles en cada anuncio.

La obtención de esta información se realizará mediante un scraper desarrollado y adaptado para este proyecto a partir del repositorio de Rocío Benítez García \cite{benitez_fotocasa_scraper}. El código original estaba orientado a la extracción de anuncios de Fotocasa para Madrid y se ha modificado progresivamente para adaptarlo al ámbito geográfico de Andalucía y a las necesidades específicas del TFM. Entre las modificaciones realizadas se encuentran la ejecución por provincias, la adaptación de las búsquedas territoriales, la deduplicación de anuncios y la incorporación de mecanismos de control y reanudación del proceso de extracción. El objetivo es recopilar de forma estructurada la información disponible en los anuncios y construir posteriormente el conjunto de datos utilizado para el análisis y el entrenamiento de los modelos.

Durante este proceso se aplicarán mecanismos de control de calidad para evitar, en la medida de lo posible, registros duplicados, anuncios incompletos o inconsistencias derivadas del proceso de extracción. También se analizará la cobertura obtenida en las distintas provincias y municipios, ya que la distribución de anuncios disponibles puede ser desigual y afectar a la representatividad territorial del conjunto de datos.

Como fuente complementaria se utilizará información oficial publicada por el Ministerio de Vivienda y Agenda Urbana (MIVAU), en particular el conjunto de datos de transacciones inmobiliarias de vivienda \cite{mivau_transacciones}. Esta fuente proporciona información agregada sobre la actividad de compraventa y, por tanto, tiene una naturaleza diferente de los microdatos obtenidos de los anuncios. Su función dentro del trabajo será principalmente contextual y, cuando el nivel territorial y temporal de agregación lo permita, podrá estudiarse su utilización para generar variables adicionales relacionadas con el mercado inmobiliario de cada zona.

Esta distinción resulta importante porque el objetivo del TFM es predecir el precio de oferta de viviendas individuales. Por ello, la variable objetivo procederá de los anuncios inmobiliarios, mientras que los datos agregados de fuentes oficiales se utilizarán como información complementaria y no como sustitutos del precio de cada inmueble.

La obtención y tratamiento de los datos se realizará teniendo en cuenta las condiciones de uso de las fuentes y limitando la recopilación a la información necesaria para los objetivos académicos del trabajo. No se almacenarán nombres, teléfonos, correos electrónicos ni otros datos personales asociados a los anunciantes. La información obtenida se utilizará exclusivamente para el análisis desarrollado en el TFM y no se plantea la redistribución del conjunto de datos bruto obtenido mediante el proceso de extracción.

\section{Métodos tradicionales de valoración y predicción}

Antes de la incorporación generalizada de técnicas de aprendizaje automático, la valoración inmobiliaria se ha abordado principalmente mediante métodos estadísticos. Entre los enfoques más utilizados se encuentran los modelos de precios hedónicos, que parten de la idea de que el precio de un bien puede explicarse a partir del conjunto de características que lo componen.

El trabajo de Rosen \cite{rosen1974hedonic} constituye una de las referencias fundamentales de este enfoque. El modelo hedónico considera que un bien diferenciado, como una vivienda, puede entenderse como una combinación de atributos y que el precio observado refleja implícitamente el valor que el mercado asigna a estas características. Aplicado al mercado inmobiliario, este planteamiento permite estudiar la relación entre el precio de una vivienda y variables como su superficie, localización, número de habitaciones o características del entorno.

En la práctica, este enfoque se ha aplicado habitualmente mediante modelos de regresión, especialmente regresión lineal. Estos modelos presentan la ventaja de ser relativamente sencillos de interpretar, ya que permiten analizar la relación entre cada variable explicativa y el precio estimado. Sin embargo, su capacidad predictiva puede verse limitada cuando existen relaciones no lineales, interacciones complejas entre variables o diferencias espaciales importantes.

Esta limitación aparece en varios de los trabajos revisados. Mora-Garcia et al. \cite{moragarcia2022housing} utilizan una regresión lineal como referencia frente a distintos algoritmos de aprendizaje automático, mientras que Guliker et al. \cite{guliker2022spatial} comparan la regresión lineal y la regresión geográficamente ponderada con XGBoost. Este tipo de comparaciones permite utilizar los modelos tradicionales como punto de referencia para determinar si técnicas más complejas proporcionan una mejora real en la capacidad predictiva.

Por este motivo, los modelos estadísticos tradicionales continúan teniendo interés en los estudios de valoración inmobiliaria, tanto por su interpretabilidad como por su utilidad como modelo base frente al que comparar técnicas de aprendizaje automático.

\section{Aprendizaje automático aplicado a la predicción de precios}

El desarrollo de técnicas de aprendizaje automático ha ampliado las posibilidades de análisis y predicción en el mercado inmobiliario. Frente a los modelos estadísticos tradicionales, estos métodos permiten modelar relaciones no lineales, interacciones entre variables y estructuras más complejas sin necesidad de especificarlas previamente de forma explícita. Por este motivo, diferentes trabajos recientes han comparado algoritmos de aprendizaje automático con modelos de regresión tradicionales para estudiar su capacidad para estimar precios de vivienda.

\subsection{Random Forest}

Random Forest es un método de aprendizaje automático basado en la combinación de múltiples árboles de decisión. El algoritmo construye diferentes árboles a partir de muestras aleatorias de los datos y utiliza subconjuntos aleatorios de variables durante su entrenamiento. En problemas de regresión, la predicción final se obtiene combinando las predicciones realizadas por los distintos árboles.

Breiman \cite{breiman2001randomforest} propuso este enfoque como una forma de mejorar la capacidad de generalización de los árboles individuales mediante la combinación de múltiples predictores. Entre sus características se encuentra la posibilidad de trabajar con relaciones no lineales y de obtener medidas de importancia de las variables, además de ser aplicable tanto a problemas de clasificación como de regresión.

En el ámbito inmobiliario, Random Forest ha sido utilizado como uno de los algoritmos de referencia en diferentes estudios de predicción de precios. Mora-Garcia et al. \cite{moragarcia2022housing}, por ejemplo, lo incluyen dentro de la comparación de modelos aplicada al mercado inmobiliario de Alicante. Su utilización en este tipo de problemas resulta adecuada por la capacidad del algoritmo para representar relaciones complejas entre características estructurales, espaciales y socioeconómicas sin asumir una relación lineal entre ellas.

No obstante, al igual que ocurre con otros modelos basados en conjuntos de árboles, la interpretación directa del modelo es más compleja que en una regresión lineal. Aunque pueden obtenerse medidas de importancia de las variables, comprender la contribución concreta de cada característica a una predicción determinada requiere técnicas adicionales de interpretabilidad. Esta limitación resulta relevante en el presente trabajo, en el que no solo se pretende comparar la capacidad predictiva de los modelos, sino también analizar los factores que influyen en sus resultados.

\subsection{Modelos de Gradient Boosting}

Los métodos de Gradient Boosting construyen un modelo predictivo mediante la combinación secuencial de múltiples modelos sencillos, generalmente árboles de decisión \cite{friedman2001gradient}. A diferencia de métodos como Random Forest, en los que los árboles se entrenan de forma más independiente, en boosting cada nuevo modelo intenta corregir los errores cometidos por los modelos anteriores.

Este enfoque permite construir modelos con una gran capacidad para representar relaciones no lineales e interacciones complejas entre variables. Por este motivo, los algoritmos basados en Gradient Boosting se han utilizado ampliamente en problemas de regresión con datos tabulares y han adquirido una relevancia importante en tareas de predicción de precios inmobiliarios.

Dentro de esta familia existen diferentes implementaciones que introducen mejoras relacionadas con la eficiencia computacional, el tratamiento de las variables o la regularización del modelo. Entre las más utilizadas se encuentran XGBoost, LightGBM y CatBoost, que presentan diferencias en su forma de construir los árboles y en el tratamiento de los datos.

En el presente trabajo se consideran estos tres algoritmos por su capacidad para trabajar con relaciones complejas entre las características de las viviendas y el precio, así como por su presencia frecuente en estudios recientes de predicción inmobiliaria. Mora-Garcia et al. \cite{moragarcia2022housing}, por ejemplo, incluyen XGBoost y LightGBM dentro de la comparación de modelos realizada sobre datos inmobiliarios del mercado de Alicante.

\subsubsection{XGBoost}

XGBoost, desarrollado por Chen y Guestrin \cite{chen2016xgboost}, es una implementación optimizada de Gradient Boosting basada en árboles de decisión. Su diseño incorpora distintas mejoras orientadas a aumentar la eficiencia computacional y controlar la complejidad del modelo, entre ellas mecanismos de regularización y estrategias optimizadas para la construcción de los árboles.

Una de las principales ventajas de XGBoost es su capacidad para trabajar con relaciones no lineales y con interacciones complejas entre variables. Esto resulta especialmente útil en problemas de predicción de precios inmobiliarios, donde el efecto de una característica puede depender de otras variables como la localización, la superficie o el tipo de vivienda.

Además, XGBoost incluye mecanismos destinados a reducir el sobreajuste y permite ajustar diferentes hiperparámetros relacionados con la profundidad de los árboles, la tasa de aprendizaje, el número de estimadores o la regularización. Esta flexibilidad permite adaptar el modelo a las características concretas del conjunto de datos utilizado.

En trabajos aplicados al mercado inmobiliario, XGBoost ha mostrado resultados competitivos frente a modelos tradicionales y otros algoritmos de aprendizaje automático. Mora-Garcia et al. \cite{moragarcia2022housing} lo incluyen entre los modelos evaluados sobre datos inmobiliarios de Alicante, mientras que Guliker et al. \cite{guliker2022spatial} lo utilizan para analizar el efecto de variables espaciales en la valoración de inmuebles.

Por estas características, XGBoost se considera uno de los algoritmos principales a evaluar en este TFM, tanto por su capacidad predictiva como por su adecuación a conjuntos de datos tabulares con variables de distinta naturaleza.

\subsubsection{LightGBM}

LightGBM es un algoritmo basado en Gradient Boosting diseñado para mejorar la eficiencia del entrenamiento cuando se trabaja con conjuntos de datos grandes o con un elevado número de variables. Ke et al. \cite{ke2017lightgbm} introdujeron este método incorporando dos técnicas principales: Gradient-based One-Side Sampling (GOSS) y Exclusive Feature Bundling (EFB).

GOSS reduce el número de observaciones utilizadas para estimar las divisiones de los árboles dando mayor importancia a aquellas que presentan gradientes elevados, mientras que EFB agrupa variables que raramente toman valores distintos de cero al mismo tiempo. Estas estrategias permiten reducir el coste computacional del entrenamiento manteniendo una capacidad predictiva competitiva.

Entre sus ventajas se encuentran la velocidad de entrenamiento, el menor consumo de memoria y la capacidad para trabajar con grandes conjuntos de datos. Estas características hacen que LightGBM sea especialmente interesante cuando se comparan varios modelos o se realizan procesos de ajuste de hiperparámetros que requieren entrenar numerosas configuraciones.

En el ámbito de la predicción inmobiliaria, Mora-Garcia et al. \cite{moragarcia2022housing} incluyen LightGBM entre los modelos evaluados sobre datos del mercado de Alicante. Su utilización permite comparar su rendimiento con otros métodos basados en árboles y estudiar si sus características de eficiencia se traducen también en una buena capacidad predictiva sobre datos inmobiliarios.

Por estas razones, LightGBM se incluye entre los algoritmos que se evaluarán en este trabajo junto con XGBoost y otros modelos de aprendizaje automático.

\subsubsection{CatBoost}

CatBoost es un algoritmo de Gradient Boosting desarrollado especialmente para trabajar de forma eficiente con variables categóricas. Prokhorenkova et al. \cite{prokhorenkova2018catboost} propusieron este método con el objetivo de reducir algunos de los problemas que pueden aparecer al codificar manualmente este tipo de variables antes del entrenamiento.

Una de sus principales características es que puede tratar variables categóricas de forma nativa mediante técnicas específicas de transformación interna. Esto puede resultar especialmente útil en conjuntos de datos inmobiliarios, donde es habitual encontrar variables como el tipo de vivienda, el municipio, la provincia, el estado del inmueble o determinadas características del anuncio.

CatBoost incorpora además una estrategia denominada \textit{ordered boosting} y procedimientos específicos para transformar las variables categóricas durante el entrenamiento. Estas técnicas fueron diseñadas para reducir el sesgo asociado al uso de información de la variable objetivo durante la codificación de categorías y limitar el denominado \textit{prediction shift} \cite{prokhorenkova2018catboost}.

Estas características hacen que CatBoost resulte especialmente interesante para el presente trabajo, ya que parte de la información disponible en los anuncios inmobiliarios tendrá previsiblemente una naturaleza categórica. Su inclusión permitirá comparar si el tratamiento nativo de este tipo de variables aporta ventajas frente a otros algoritmos como XGBoost o LightGBM.

\subsection{Redes neuronales}

Las redes neuronales constituyen otra familia de modelos capaces de representar relaciones no lineales entre las variables de entrada y la variable objetivo. Su estructura está formada por diferentes capas de unidades conectadas entre sí, que transforman progresivamente la información hasta obtener una predicción final.

Uno de los modelos más habituales para trabajar con datos tabulares es el perceptrón multicapa o MLP. Este tipo de red está compuesto por una capa de entrada, una o varias capas ocultas y una capa de salida. Durante el entrenamiento, los pesos de las conexiones se ajustan para reducir el error cometido por el modelo \cite{goodfellow2016deep}.

La principal ventaja de las redes neuronales es su capacidad para aproximar relaciones complejas entre las variables. Sin embargo, su rendimiento depende en gran medida de aspectos como la arquitectura utilizada, el número de capas y neuronas, las funciones de activación, la regularización y el proceso de entrenamiento. Además, suelen requerir un preprocesamiento más cuidadoso de los datos y pueden ser más sensibles a la escala de las variables que los modelos basados en árboles.

En el ámbito de la predicción inmobiliaria, las redes neuronales han sido utilizadas como alternativa a los modelos estadísticos y a los métodos basados en árboles. La revisión de Geerts et al. \cite{houseprice_survey2023} recoge diferentes trabajos en los que se utilizan redes neuronales para estimar precios de vivienda, lo que permite considerar este enfoque como una alternativa adicional dentro de la comparación de modelos.

En este trabajo se valorará la utilización de una red neuronal tipo MLP como modelo complementario. Su inclusión permitirá comparar un enfoque basado en redes neuronales con los modelos de árboles considerados anteriormente y analizar si aporta mejoras en la estimación del precio de oferta de las viviendas.

\section{Preparación de los datos e ingeniería de características}

La preparación de los datos constituye una fase fundamental antes del entrenamiento de los modelos predictivos. La calidad del conjunto de datos, el tratamiento de los valores ausentes y la transformación de las variables pueden influir de forma importante en el rendimiento final de los modelos.

En los problemas de predicción inmobiliaria es habitual trabajar con variables de distinta naturaleza. Pueden encontrarse variables numéricas, como la superficie o el número de habitaciones; variables categóricas, como el tipo de vivienda o el municipio; y variables relacionadas con la localización y el entorno. La revisión de Geerts et al. \cite{houseprice_survey2023} muestra precisamente la variedad de tipos de datos utilizados en este tipo de estudios.

Entre las tareas habituales de preprocesamiento se encuentran la detección de valores ausentes, la identificación de valores atípicos, la eliminación de registros duplicados y la revisión de posibles inconsistencias. Estas operaciones permiten reducir errores en los datos y evitar que determinados registros puedan distorsionar el proceso de aprendizaje.

También puede ser necesario transformar algunas variables antes de utilizarlas en los modelos. Las variables categóricas pueden requerir técnicas de codificación, mientras que determinadas variables numéricas pueden beneficiarse de procesos de escalado o normalización, especialmente en modelos sensibles a la escala como las redes neuronales. En cambio, los modelos basados en árboles suelen ser menos sensibles a este tipo de transformaciones.

La ingeniería de características permite además crear nuevas variables a partir de la información disponible. En el contexto inmobiliario, pueden resultar útiles variables derivadas como indicadores relacionados con la antigüedad de la vivienda, agrupaciones territoriales, características obtenidas a partir de otras variables del inmueble o medidas relacionadas con la proximidad a determinados servicios. La utilidad de estas variables deberá evaluarse evitando incorporar información derivada directa o indirectamente de la variable objetivo, ya que esto podría producir fugas de información y generar estimaciones artificialmente optimistas del rendimiento de los modelos.

La selección definitiva de las técnicas de preprocesamiento e ingeniería de características dependerá de la estructura y calidad del conjunto de datos obtenido. En este trabajo se analizarán estas decisiones antes del entrenamiento de los modelos, procurando aplicar el mismo criterio de preparación cuando sea necesario para garantizar una comparación adecuada entre ellos.

\section{Evaluación de los modelos predictivos}

La evaluación de los modelos es necesaria para comparar de forma objetiva su capacidad para estimar el precio de las viviendas. En problemas de regresión pueden utilizarse diferentes métricas, cada una de las cuales permite analizar el error desde una perspectiva distinta.

Una de las métricas más utilizadas es el error absoluto medio o MAE, que calcula la media de las diferencias absolutas entre los valores reales y las predicciones realizadas por el modelo. Su principal ventaja es que se expresa en las mismas unidades que la variable objetivo, lo que facilita la interpretación del error medio cometido.

Otra métrica habitual es el error cuadrático medio, normalmente expresado mediante su raíz cuadrada o RMSE. Esta medida penaliza con mayor intensidad los errores grandes, ya que las diferencias entre los valores reales y predichos se elevan al cuadrado antes de calcular la media. Por este motivo, puede resultar especialmente útil cuando se quiere tener en cuenta el impacto de predicciones con desviaciones elevadas.

También puede utilizarse el coeficiente de determinación $R^2$, que permite analizar qué proporción de la variabilidad de la variable objetivo es explicada por el modelo. A diferencia de MAE y RMSE, donde valores más bajos indican un menor error, en $R^2$ los valores más próximos a 1 representan generalmente una mayor capacidad explicativa.

La utilización conjunta de estas métricas permite obtener una visión más completa del comportamiento de los modelos. Un modelo puede presentar un buen error medio y, al mismo tiempo, cometer errores elevados en determinados casos, por lo que resulta conveniente analizar más de una medida antes de seleccionar la alternativa con mejor rendimiento.

En el presente trabajo se utilizarán principalmente MAE, RMSE y $R^2$ para comparar los distintos modelos entrenados. La evaluación se realizará sobre los mismos conjuntos de datos y siguiendo un procedimiento común, con el objetivo de que las diferencias observadas puedan atribuirse al comportamiento de los modelos y no a cambios en las condiciones de evaluación.

\subsection{Estrategia de validación}

Además de seleccionar métricas adecuadas, será necesario definir un procedimiento de validación que permita estimar de forma realista la capacidad de generalización de los modelos. Para ello, el conjunto de datos se dividirá en subconjuntos independientes de entrenamiento, validación y prueba. Los datos de entrenamiento se utilizarán para ajustar los modelos, los datos de validación servirán para comparar configuraciones y realizar el ajuste de hiperparámetros, y el conjunto de prueba se reservará para la evaluación final de los modelos seleccionados.

El conjunto de prueba no se utilizará durante la selección de variables, el ajuste de hiperparámetros ni la elección del modelo final. Del mismo modo, cualquier transformación que dependa de la información contenida en los datos, como la imputación de valores ausentes, el escalado o determinadas codificaciones de variables, se ajustará utilizando únicamente los datos de entrenamiento y se aplicará posteriormente a los conjuntos de validación y prueba. De esta forma se reducirá el riesgo de fuga de información entre las diferentes etapas del proceso.

En el caso de los datos procedentes de anuncios inmobiliarios será especialmente importante detectar registros duplicados o posibles republicaciones de una misma vivienda antes de realizar la división de los datos. La presencia de observaciones correspondientes al mismo inmueble, o a anuncios prácticamente idénticos, tanto en entrenamiento como en prueba podría producir una estimación artificialmente optimista del rendimiento. Por este motivo, la deduplicación formará parte de la preparación previa del conjunto de datos.

También se tendrá en cuenta la componente territorial del problema. Una división aleatoria permite comparar los modelos bajo unas mismas condiciones, pero puede no reflejar completamente su comportamiento cuando se aplican a zonas geográficas diferentes de las utilizadas durante el entrenamiento. Por ello, siempre que la cantidad y distribución de los datos lo permitan, la evaluación general se complementará con un análisis del rendimiento por provincias o ámbitos territoriales. También se valorará realizar una validación adicional mediante agrupaciones geográficas para estudiar hasta qué punto los modelos mantienen su capacidad predictiva al cambiar de zona.

La estrategia definitiva de validación se determinará una vez conocida la cobertura y distribución del conjunto de datos obtenido. En cualquier caso, se mantendrá un procedimiento común para todos los modelos y se documentarán las decisiones adoptadas para que las diferencias de rendimiento puedan interpretarse de forma comparable.

\section{Optimización de hiperparámetros}

El rendimiento de muchos modelos de aprendizaje automático depende de la configuración de sus hiperparámetros. Estos parámetros no se aprenden directamente durante el entrenamiento, sino que deben definirse previamente y pueden influir de forma importante en la capacidad predictiva, la complejidad y el riesgo de sobreajuste del modelo.

Entre los hiperparámetros más habituales se encuentran la profundidad de los árboles, la tasa de aprendizaje, el número de estimadores, los parámetros de regularización o, en el caso de las redes neuronales, el número de capas, el número de neuronas y la tasa de aprendizaje utilizada durante el entrenamiento.

Una estrategia sencilla consiste en probar manualmente diferentes configuraciones, aunque este procedimiento puede resultar poco eficiente cuando el número de hiperparámetros es elevado. También pueden utilizarse técnicas como la búsqueda en rejilla o Grid Search y la búsqueda aleatoria, que permiten explorar automáticamente distintas combinaciones de valores.

En este trabajo se plantea utilizar Optuna \cite{akiba2019optuna} como herramienta de optimización de hiperparámetros. Optuna permite definir un espacio de búsqueda y explorar diferentes configuraciones del modelo de forma automática, utilizando los resultados obtenidos en evaluaciones anteriores para orientar la búsqueda hacia combinaciones potencialmente más adecuadas.

La optimización se aplicará únicamente cuando resulte necesaria y se realizará utilizando los datos de entrenamiento y validación, evitando utilizar el conjunto de prueba durante este proceso. De esta forma, el conjunto de prueba podrá reservarse para realizar una evaluación final del rendimiento de los modelos seleccionados.

\section{Interpretabilidad de los modelos}

La interpretabilidad de los modelos es un aspecto importante cuando se utilizan algoritmos de aprendizaje automático con una estructura compleja. Aunque modelos como Random Forest, XGBoost, LightGBM, CatBoost o las redes neuronales pueden ofrecer una buena capacidad predictiva, su funcionamiento interno resulta menos transparente que el de modelos más sencillos como una regresión lineal.

En este contexto, las técnicas de interpretabilidad permiten analizar qué variables influyen en las predicciones realizadas por el modelo y en qué dirección lo hacen. Esto resulta especialmente relevante en problemas como la predicción de precios inmobiliarios, donde no solo interesa obtener una estimación precisa, sino también comprender qué características de la vivienda o de su entorno están asociadas a aumentos o disminuciones en el precio estimado.

Una de las técnicas más utilizadas para este propósito es SHAP, propuesta por Lundberg y Lee \cite{lundberg2017shap}. Este método se basa en conceptos procedentes de la teoría de juegos y asigna a cada variable una contribución concreta respecto a la predicción realizada por el modelo.

La aplicación de técnicas de explicabilidad en el ámbito inmobiliario aparece también en trabajos recientes. Iban \cite{iban2022explainable} utiliza SHAP para interpretar modelos basados en árboles empleados en valoración inmobiliaria, incluyendo Random Forest, XGBoost y LightGBM. Por su parte, Teoh et al. \cite{teoh2023explainable} combinan modelos predictivos con análisis SHAP para estudiar qué características presentan una mayor influencia en la estimación de precios de vivienda. Estos trabajos muestran que la interpretabilidad puede utilizarse como complemento de las métricas predictivas para comprender mejor el comportamiento de los modelos.

Los valores SHAP pueden utilizarse tanto para realizar interpretaciones globales como locales. A nivel global permiten analizar qué variables tienen una mayor influencia sobre el conjunto de predicciones, mientras que a nivel local permiten estudiar cómo contribuye cada característica a la predicción de una vivienda concreta.

A pesar de su utilidad, los valores SHAP deben interpretarse con cautela. Estas explicaciones describen cómo utiliza el modelo las variables para construir sus predicciones, pero no permiten establecer relaciones causales entre una característica y el precio de la vivienda. Además, cuando existen variables fuertemente relacionadas entre sí, la distribución de la importancia entre ellas puede dificultar la interpretación individual de sus contribuciones.

En el presente trabajo se utilizará SHAP para complementar la evaluación predictiva de los modelos seleccionados. El objetivo será identificar las variables que presentan una mayor influencia en las estimaciones de precio y analizar si estos resultados son coherentes con las relaciones observadas durante el análisis exploratorio y con lo descrito en la literatura revisada.

\section{Antecedentes en el mercado inmobiliario español y andaluz}

En España existen distintos trabajos que han aplicado técnicas estadísticas y de aprendizaje automático al análisis del mercado inmobiliario. Entre ellos, el estudio de Mora-Garcia et al. \cite{moragarcia2022housing} resulta especialmente relevante por analizar datos inmobiliarios georreferenciados de Alicante y comparar una regresión lineal con diferentes algoritmos de aprendizaje automático. Los resultados muestran un mejor comportamiento de los modelos no lineales, especialmente de los métodos basados en boosting, con valores de $R^2$ próximos a 0,92 en los modelos con mejor rendimiento. Los autores señalan además que algunos algoritmos basados en bagging, como Random Forest y Extra Trees, presentan mayores problemas de sobreajuste, mientras que XGBoost y LightGBM ofrecen un mejor equilibrio entre capacidad predictiva, generalización y coste computacional.

También en Alicante, Rico-Juan y Taltavull de La Paz \cite{ricojuan2021asking} analizan específicamente los precios de oferta del mercado residencial mediante modelos hedónicos y técnicas de aprendizaje automático. En este caso, Random Forest obtiene el mejor rendimiento predictivo entre los enfoques comparados y los autores señalan que permite estimar los valores inmobiliarios con un error inferior al 2\%. Este resultado muestra el potencial de los modelos basados en árboles para trabajar con precios de oferta, aunque sus conclusiones se encuentran condicionadas por el ámbito territorial concreto de Alicante y no pueden trasladarse directamente a otros mercados inmobiliarios.

En el caso de Andalucía, Gutiérrez et al. \cite{gutierrez2026andalusia} estudian la predicción de precios de vivienda en seis ciudades andaluzas utilizando datos procedentes de anuncios inmobiliarios en línea. El conjunto de datos combina características estructurales de las viviendas con variables socioeconómicas y de localización. Los autores comparan un modelo hedónico estimado mediante mínimos cuadrados ordinarios con Random Forest, K-Nearest Neighbours y XGBoost. Los resultados muestran un peor rendimiento de los modelos lineales frente a los métodos de aprendizaje automático, lo que pone de manifiesto la dificultad de representar mediante una única relación lineal la heterogeneidad del mercado inmobiliario.

Además del rendimiento predictivo, los autores aplican valores SHAP al modelo XGBoost para identificar las variables con mayor influencia en cada ciudad. Los resultados muestran que las características estructurales mantienen una importancia elevada, mientras que el efecto de las variables socioeconómicas y de localización varía entre ciudades \cite{gutierrez2026andalusia}. Esta variabilidad constituye un aspecto especialmente relevante para un análisis a escala andaluza, ya que indica que las relaciones aprendidas en una ciudad no tienen por qué reproducirse de la misma forma en otros mercados locales. No obstante, aunque el estudio incorpora seis ciudades andaluzas, su ámbito no cubre de forma completa el territorio de Andalucía, por lo que sus resultados deben interpretarse teniendo en cuenta esta limitación geográfica y la posible heterogeneidad existente entre municipios no incluidos en el análisis.

Otro antecedente directamente relacionado con este TFM es el estudio de Gutiérrez et al. \cite{gutierrez2026sevilla} centrado en la ciudad de Sevilla. El trabajo utiliza 1.710 anuncios inmobiliarios e incorpora características estructurales, indicadores socioeconómicos a nivel de sección censal y variables de proximidad a servicios. Los autores comparan una regresión hedónica mediante mínimos cuadrados ordinarios con un modelo Random Forest y obtienen un mejor rendimiento predictivo con este último, reduciendo tanto el MAE como el RMSE respecto al modelo lineal.

El estudio de Sevilla incorpora además valores SHAP, curvas ICE y gráficos PDP para interpretar las predicciones. Sus resultados muestran que variables socioeconómicas como la renta media y la proporción de población con estudios superiores, junto con determinadas características estructurales de la vivienda, presentan una influencia importante en los precios de oferta \cite{gutierrez2026sevilla}. No obstante, el análisis se limita a una única ciudad y a una muestra de 1.710 anuncios, por lo que sus resultados no permiten asumir que las mismas relaciones se mantengan de forma homogénea en el conjunto de Andalucía.

Estos antecedentes presentan una relación directa con el enfoque planteado en este TFM. En particular, respaldan la utilización de anuncios inmobiliarios como fuente de información, la incorporación de variables relacionadas con el entorno, la comparación entre modelos tradicionales y métodos de aprendizaje automático y el uso de técnicas de interpretabilidad. El presente trabajo ampliará este planteamiento mediante la comparación de varios algoritmos, incluyendo Random Forest, XGBoost, LightGBM, CatBoost y, si los datos lo permiten, una red neuronal tipo MLP, manteniendo como ámbito de estudio el mercado residencial andaluz.

\subsection{Hueco de investigación y aportación del trabajo}

Los antecedentes revisados muestran que ya existen estudios que aplican técnicas de aprendizaje automático a la predicción de precios inmobiliarios en España y, de forma más específica, en Andalucía. Sin embargo, su alcance, los modelos comparados y la forma de abordar la interpretabilidad no coinciden completamente con el planteamiento de este TFM. El estudio de Gutiérrez et al. \cite{gutierrez2026andalusia}, por ejemplo, analiza seis ciudades andaluzas y compara un modelo hedónico con Random Forest, K-Nearest Neighbours y XGBoost, mientras que el estudio posterior centrado en Sevilla \cite{gutierrez2026sevilla} limita el análisis a una única ciudad y compara principalmente una regresión hedónica con Random Forest.

A partir de estos antecedentes, el presente trabajo no pretende proponer un nuevo algoritmo de aprendizaje automático, sino ampliar la comparación de enfoques aplicados al mercado residencial andaluz. Para ello se plantea evaluar, bajo un mismo procedimiento de preparación y evaluación de los datos, una regresión lineal utilizada como referencia, Random Forest y varios algoritmos de Gradient Boosting, concretamente XGBoost, LightGBM y CatBoost. También se valorará la utilización de una red neuronal tipo MLP como modelo complementario cuando las características del conjunto de datos lo permitan.

Otra diferencia del trabajo será el análisis conjunto de la capacidad predictiva y de la interpretabilidad de los modelos. Además de comparar su rendimiento mediante métricas comunes, se utilizará SHAP para estudiar la influencia de las variables sobre las predicciones y analizar si los factores identificados son coherentes entre distintos modelos y ámbitos territoriales. Asimismo, se estudiará la posibilidad de enriquecer la información procedente de los anuncios inmobiliarios con variables obtenidas de fuentes oficiales cuando puedan integrarse de forma adecuada.

De este modo, la aportación del TFM se sitúa principalmente en la aplicación comparativa y sistemática de un conjunto más amplio de modelos sobre el mercado residencial andaluz, manteniendo como variable objetivo el precio de oferta y prestando especial atención tanto al rendimiento predictivo como a la interpretación de los resultados. El alcance territorial definitivo y las fuentes adicionales utilizadas dependerán de la cobertura y calidad de los datos finalmente disponibles.

El Cuadro \ref{tab:literatura_inmobiliaria} resume algunos de los trabajos más relacionados con el enfoque de este TFM y permite comparar su ámbito de estudio, la variable objetivo, los principales modelos utilizados y la incorporación de técnicas de interpretabilidad.

\begin{table}[H]
\centering
\caption{Comparación de trabajos relevantes sobre predicción y valoración inmobiliaria.}
\label{tab:literatura_inmobiliaria}
\small
\renewcommand{\arraystretch}{1.35}
\begin{tabularx}{\textwidth}{
>{\raggedright\arraybackslash}p{4.3cm}
>{\raggedright\arraybackslash}X
}

\toprule
\textbf{Trabajo} &
\textbf{Principales características}
\tabularnewline
\midrule
Mora-Garcia et al. \cite{moragarcia2022housing} &
Alicante. Predicción de precios de vivienda mediante regresión lineal, Random Forest, Extra Trees, Gradient Boosting, XGBoost y LightGBM. Incluye análisis de importancia de variables.
\tabularnewline
\midrule
Rico-Juan y Taltavull de La Paz \cite{ricojuan2021asking} &
Alicante. Predicción del precio de oferta mediante modelos hedónicos y técnicas de aprendizaje automático, incorporando además análisis de explicabilidad.
\tabularnewline
\midrule
Guliker et al. \cite{guliker2022spatial} &
Países Bajos. Valoración inmobiliaria mediante regresión lineal, regresión geográficamente ponderada y XGBoost, con especial atención a las variables espaciales.
\tabularnewline
\midrule
Iban \cite{iban2022explainable} &
Valoración residencial mediante Random Forest, XGBoost, LightGBM y Gradient Boosting. Utiliza SHAP para interpretar los factores que influyen en las estimaciones.
\tabularnewline
\midrule
Gutiérrez et al. \cite{gutierrez2026andalusia} &
Seis ciudades de Andalucía. Predicción utilizando anuncios inmobiliarios y comparación de regresión hedónica, Random Forest, KNN y XGBoost. Incorpora SHAP.
\tabularnewline
\midrule
Gutiérrez et al. \cite{gutierrez2026sevilla} &
Sevilla. Análisis del precio de oferta mediante regresión hedónica y Random Forest, utilizando técnicas de interpretabilidad como SHAP.
\tabularnewline
\bottomrule
\end{tabularx}
\end{table}

La comparación muestra que los trabajos más próximos al presente TFM combinan progresivamente modelos tradicionales y técnicas de aprendizaje automático, incorporando además variables espaciales, socioeconómicas y herramientas de interpretabilidad. También se observa que no todos los estudios utilizan la misma variable objetivo, ya que algunos trabajan con precios de oferta, otros con precios inmobiliarios obtenidos de anuncios y otros con valores de tasación. Esta diferencia refuerza la necesidad de interpretar los resultados de cada estudio dentro de su propio contexto.

\section{Síntesis del estado del arte}

La revisión realizada muestra que la predicción de precios de vivienda se ha abordado mediante enfoques muy diversos, desde modelos estadísticos tradicionales hasta algoritmos de aprendizaje automático capaces de representar relaciones no lineales e interacciones complejas entre variables. Los modelos hedónicos y la regresión lineal continúan siendo útiles como referencia por su facilidad de interpretación, mientras que métodos como Random Forest y los algoritmos de Gradient Boosting ofrecen una mayor flexibilidad para trabajar con relaciones más complejas.

La literatura revisada también pone de manifiesto la importancia de la calidad y variedad de la información utilizada. Además de las características estructurales de las viviendas, las variables relacionadas con la localización, el entorno urbano y las condiciones socioeconómicas pueden aportar información relevante para explicar las diferencias de precio entre inmuebles y zonas geográficas.

Los trabajos desarrollados en España y, especialmente, los centrados en Andalucía muestran que la aplicación de técnicas de aprendizaje automático al mercado inmobiliario regional es viable y que los modelos no lineales pueden ofrecer resultados competitivos frente a enfoques tradicionales. Estos antecedentes también respaldan la utilización de datos procedentes de anuncios inmobiliarios y la incorporación de técnicas de interpretabilidad para analizar los factores que influyen en las predicciones.

La revisión realizada permite además concretar algunos de los objetivos parciales planteados inicialmente. En particular, la comparación de modelos se centrará prioritariamente en una regresión lineal utilizada como referencia, Random Forest y distintos algoritmos de Gradient Boosting como XGBoost, LightGBM y CatBoost. La utilización de una red neuronal MLP se mantendrá como alternativa complementaria en función de las características y del volumen de los datos disponibles. Asimismo, la evaluación se apoyará principalmente en MAE, RMSE y $R^2$, se utilizará Optuna para el ajuste de hiperparámetros cuando resulte necesario y SHAP para estudiar la contribución de las variables a las predicciones.

Además de analizar la capacidad predictiva de los modelos, se prestará especial atención a su interpretación mediante SHAP. Esto permitirá estudiar qué variables tienen una mayor influencia sobre las estimaciones realizadas y analizar si estas relaciones son coherentes con los resultados descritos en la literatura.

En conjunto, el estado del arte proporciona una base para la metodología planteada en este TFM y permite justificar tanto la selección inicial de modelos como las técnicas de evaluación, optimización e interpretabilidad que se utilizarán posteriormente. La implementación final podrá ajustarse en función de la calidad y características del conjunto de datos obtenido, manteniendo como objetivo principal la estimación del precio de oferta de viviendas residenciales en Andalucía.

\iffalse
\chapter{Materiales y métodos}
En estos apartados, hay que describir:
\begin{itemize}
    \item Los aspectos más relevante del diseño y desarrollo del trabajo.
    \item La metodología elegida para hacer este desarrollo, describiendo las alternativas posibles, las decisiones tomadas, y los criterios utilizados para tomar estas decisiones.
    \item Los productos obtenidos.
\end{itemize}

 
\textbf{La estructuración de los capítulos puede variar según el tipo de trabajo.}  
 
En caso de que se proceda, se incluirá un apartado de “Valoración económica del trabajo”. Este apartado indicará los gastos asociados al desarrollo y mantenimiento del trabajo, así como los beneficios económicos obtenidos y un análisis final sobre la viabilidad del producto.



\chapter{Resultados}

Detallad en este apartado los resultados obtenidos utilizando la metodología descrita en el apartado anterior.

Las figuras tienen que estar explicadas y citadas en el texto, como la \ref{fig:my_label}, en la cual se muestra el error en función de la distancia, en unidades arbitrarias. En todas las gráficas tiene que haber el título de los ejes.

\begin{figure}[!htbp]
    \centering
    \includegraphics[width=7truecm]{Rplotmanh.png}
    \caption{Error en función de la distancia en unidades arbitrarias.}
    \label{fig:my_label}
\end{figure}

%\chapter{Discusión}
%Discusión de los resultados en el contexto del proyecto. Es en este apartado donde cobran sentido y en el cual se responden las preguntas de investigación y se muestra como los resultados dan respuesta a los problemas planteados.

%Esta parte puede ser que no aplique según el tipo de trabajo.

%\chapter{Valoración económica}
%En caso de que corresponda, se incluirá un apartado de ``Valoración económica del trabajo". Este apartado indicará los gastos asociados al desarrollo y mantenimiento del trabajo, así como los beneficios económicos obtenidos. Hay que hacer un análisis final sobre la viabilidad del producto.

\chapter{Conclusiones y trabajos futuros}

%\section{Conclusiones}
Este capítulo tiene que incluir:
\begin{itemize}
\item Una descripción de las conclusiones del trabajo:
\begin{itemize}
    \item Una vez se han obtenido los resultados, ¿qué conclusiones se extraen?
    \item ¿Estos resultados son los esperados? ¿O han sido sorprendentes? ¿Por qué? 
\end{itemize}
\item Una reflexión crítica sobre el logro de los objetivos planteados inicialmente:
\begin{itemize}
    \item ¿Hemos logrado todos los objetivos? Si la respuesta es negativa, ¿por qué motivo?
\end{itemize}
\item Un análisis crítico del seguimiento de la planificación y metodología a lo largo del producto:
\begin{itemize}
    \item ¿Se ha seguido la planificación?
    \item ¿La metodología prevista ha sido suficientemente adecuada?
    \item ¿Ha habido que introducir cambios para garantizar el éxito del trabajo? ¿Por qué? 
\end{itemize}
\item De los impactos previstos en \ref{s:etic}, ético-sociales, de sostenibilidad y de diversidad, evaluar/mencionar si se han mitigado (si eran negativos) o si se han conseguido (si eran positivos). 
\item Si han aparecido impactos no previstos a \ref{s:etic}, evaluar/mencionar cómo se han mitigado (si eran negativos) o que han aportado (si eran positivos).
\item Las líneas de trabajo futuro que no se han podido explorar en este trabajo y han quedado pendientes.


%\item Una descripción de las conclusiones del trabajo: Qué lecciones se han aprendido del trabajo?
%\item Una reflexión crítica sobre el logro de los objetivos planteados inicialmente: Hemos conseguido todos los objetivos? Si la respuesta es negativa, por qué motivo?
%\item De los impactos previstos a la sección \ref{s:etic}, una evaluación, o al menos, mención, sobre si se han mitigado (si eran negativos) o si se han conseguido (si eran positivos). 
\end{itemize}

%\section{Líneas de futuro}
%Las líneas de trabajo futuro que no se han podido explorar en este trabajo y han quedado pendientes.

%\section{Seguimiento de la planificación}
%Un análisis crítico del seguimiento de la planificación y metodología a lo largo del trabajo: 
%\begin{itemize}
%    \item  Se ha seguido la planificación? 
%    \item La metodología prevista ha sido la adecuada? 
%    \item Ha habido que introducir cambios para garantizar el éxito del trabajo? ¿Por qué? 
%\end{itemize}

\chapter{Glosario}

Definición de los términos y acrónimos más relevantes utilizados dentro de la Memoria.

\fi

\bibliographystyle{plain}
\bibliography{bibliografia}

\end{document}
